\chapter{GUIDANCE SEQUENCE TESTS AND EVALUATION RUBRIC}
\label{app:scenarios}
\newpage
\section{Controlled Scenario Matrix}
Each scenario records input observations, policy and configuration, eligible actions, provider behavior, playback occupancy, and the speech actually admitted. Every policy is tested on the same trajectories. The expected behaviors in the table are acceptance criteria, not reports of wearer experience.

\begingroup
\thesislongtablecaption{Sequence scenarios and acceptance criteria.}
\begin{longtable}{@{}p{0.25\linewidth}p{0.66\linewidth}@{}}
\toprule
Scenario & Evidence and expected behavior \\\midrule\endfirsthead
\toprule Scenario & Evidence and expected behavior \\\midrule\endhead
Large turn & Target opposite starting reference; progress cues remain directionally valid as plane boundaries are crossed.\\
Plane entry & Stable entry changes to refinement; looking upward on the same object plane does not create a second arrival.\\
Vertical refinement & Target above/below gaze reference; magnitude wording follows documented plane-relative thresholds.\\
Overshoot and return & A valid recent crossing supports a return cue; an outdated target bearing alone does not.\\
Gaze only & Winning vicinity qualifies with controller elsewhere; gaze acknowledgment does not say the controller is on target.\\
Controller entry & Actual target intersection suppresses contradictory directions before hold qualification.\\
Controller hold & Qualified hold plus current valid nearby gaze admits pointing acknowledgment; no trigger completion is inferred.\\
Controller departure & Leaving before qualification cancels pointing acknowledgment and allows appropriate guidance after current-state checks.\\
Reacquisition & Valid gaze away rearms confirmation; continuously holding gaze does not repeat it indefinitely.\\
Overlapping regions & Only the closest eligible vicinity candidate receives acknowledgment; neighbor overlap alone is insufficient.\\
Blink or stale gaze & Missing evidence counts as neither a gaze hold nor a departure; new gaze-dependent cues are suspended.\\
Pause and recovery & Pending cues do not carry over across a reset or pause; valid recovery resumes from fresh evidence.\\
Late model response & Previously valid direction is discarded if current controller/gaze state conflicts.\\
Provider wait/failure & Deferral remains bounded; fallback respects current eligibility and repetition limits.\\
Speech occupied & Higher-priority event follows explicit interruption/admission policy; interrupted text is not recorded as fully heard.\\
Correct trigger & Trial ends only on correct selection; success feedback uses completed-trial context.\\\bottomrule
\end{longtable}
\endgroup

\section{Deterministic Labels}
A wrong-direction label requires supported geometry showing that the delivered cue moves the user away from the intended target relation under the cue's reference frame. A stale label requires evidence that the relevant state changed before playback. Contradictory ordering requires a delivered sequence inconsistent with valid current evidence, such as an away-direction during a target controller hit. When the trace lacks the needed observation, the label is unknown.

Confirmation metrics distinguish gaze acknowledgment, controller acknowledgment, and trial completion. A gaze episode begins at its first valid qualification and is not restarted for each cloud request; a controller episode uses its own hold and gaze requirements, while completion requires correct trigger selection. The evaluator records missed acknowledgment, delayed acknowledgment, and disqualification separately rather than classifying every silent interval as a failure.

\section{Advisory Model Review}
The offline judge receives a bounded sequence containing current and recent state, candidate and chosen actions, actual playback, timing, and provenance. It evaluates whether the wording was useful, whether timing supported the task, whether repetition was unnecessary, and whether interruption preserved intelligibility. The judge is instructed to cite the concrete events that support each judgment and to distinguish uncertainty from a negative rating.

Logs and participant text are supplied to the judge as data, not as instructions. Its output is admitted to the viewer only if it matches the requested case and schema. Invalid or failed evaluations remain visibly failed or pending rather than defaulting to a favorable rating. When human reviewers disagree with the judge, adjudication preserves both the original rating and the reason for revision.

\section{Reporting Template}
Each comparison reports the recording/scenario manifest, policy revisions, number of trials, unique episodes, provider calls, latency distributions, unknown evidence, and remaining failures, with representative timelines showing qualification, request, response, cancellation, and actual playback. Simulated scheduling time, desktop network response time, headset software playback time, and measured acoustic latency are reported separately. Replay results characterize system behavior and do not demonstrate changes in participant behavior.
